\section{模型的建立与求解}
\subsection{数据预处理}
说明数据来源、清洗规则、缺失值与异常值处理，以及后续计算使用的数据形式。数据统计以实际处理结果为准。

\subsection{问题一的模型建立与求解}
\subsubsection{模型建立}
在此给出变量、目标函数、约束条件和推导过程。下面的均值公式仅用于演示公式排版，并非针对赛题提出的模型：
\begin{equation}
  \bar{x}=\frac{1}{n}\sum_{i=1}^{n}x_i .
  \label{eq:example}
\end{equation}
公式\eqref{eq:example}中的符号应在正文或符号表中解释。

\subsubsection{求解方法}
在此说明算法步骤、参数设置、停止条件与计算环境。必要时说明数值稳定性、复杂度或可复现的方法。

\subsubsection{结果与分析}
在此填写真实求解结果，并通过图表进行分析。下表只有占位文字，不代表实验结果。
\begin{table}[htbp]
  \centering
  \caption{结果展示格式示例}
  \label{tab:results}
  \begin{tabularx}{\textwidth}{lXXX}
    \toprule
    方案 & 指标一 & 指标二 & 说明 \\
    \midrule
    方案一 & 待计算 & 待计算 & 待填写 \\
    方案二 & 待计算 & 待计算 & 待填写 \\
    \bottomrule
  \end{tabularx}
\end{table}

\begin{figure}[htbp]
  \centering
  \fbox{\parbox[c][35mm][c]{0.8\textwidth}{\centering 在此插入实际结果图}}
  \caption{结果图排版示例}
  \label{fig:example}
\end{figure}
图\ref{fig:example}的图注位于图下方，表\ref{tab:results}的表题位于表上方。

\subsection{问题二的模型建立与求解}
\subsubsection{模型建立}
在此说明第二个问题的新增变量、约束及模型。
\subsubsection{求解与结果分析}
在此填写求解方法、真实结果与解释。

\subsection{问题三的模型建立与求解}
按实际问题补充模型、求解过程与结果。赛题包含更多问题时，保持同样的标题层级。
